\begin{lemma}
\label{lemma-coherent-projective}
Let $R$ be a Noetherian ring.
Let $n \geq 0$ be an integer.
For every coherent sheaf $\mathcal{F}$ on $\mathbf{P}^n_R$
we have the following:
\begin{enumerate}
\item There exists an $r \geq 0$ and
$d_1, \ldots, d_r \in \mathbf{Z}$ and a surjection
$$
\bigoplus\nolimits_{j = 1, \ldots, r}
\mathcal{O}_{\mathbf{P}^n_R}(d_j)
\longrightarrow
\mathcal{F}.
$$
\item We have $H^i(\mathbf{P}^n_R, \mathcal{F}) = 0$ unless
$0 \leq i \leq n$.
\item For any $i$ the cohomology group $H^i(\mathbf{P}^n_R, \mathcal{F})$
is a finite $R$-module.
\item If $i > 0$, then
$H^i(\mathbf{P}^n_R, \mathcal{F}(d)) = 0$ for all $d$ large enough.
\item For any $k \in \mathbf{Z}$ the graded $R[T_0, \ldots, T_n]$-module
$$
\bigoplus\nolimits_{d \geq k} H^0(\mathbf{P}^n_R, \mathcal{F}(d))
$$
is a finite $R[T_0, \ldots, T_n]$-module.
\end{enumerate}
\end{lemma}

\begin{proof}
We will use that $\mathcal{O}_{\mathbf{P}^n_R}(1)$ is an ample invertible
sheaf on
the scheme $\mathbf{P}^n_R$. This follows directly from the definition
since $\mathbf{P}^n_R$ covered by the standard affine opens $D_{+}(T_i)$.
Hence by
Properties, Proposition \ref{properties-proposition-characterize-ample}
every finite type quasi-coherent $\mathcal{O}_{\mathbf{P}^n_R}$-module
is a quotient of a finite direct sum of tensor powers of
$\mathcal{O}_{\mathbf{P}^n_R}(1)$. On the other hand coherent sheaves
and finite type quasi-coherent sheaves are the same thing on projective
space over $R$ by Lemma \ref{lemma-coherent-Noetherian}. Thus we see (1).

\medskip\noindent
Projective $n$-space $\mathbf{P}^n_R$ is covered by $n + 1$ affines,
namely the standard opens $D_{+}(T_i)$, $i = 0, \ldots, n$, see Constructions,
Lemma \ref{constructions-lemma-standard-covering-projective-space}.
Hence we see that for any quasi-coherent
sheaf $\mathcal{F}$ on $\mathbf{P}^n_R$
we have $H^i(\mathbf{P}^n_R, \mathcal{F}) = 0$ for $i \geq n + 1$,
see Lemma \ref{lemma-vanishing-nr-affines}. Hence (2) holds.

\medskip\noindent
Let us prove (3) and (4) simultaneously for all coherent sheaves
on $\mathbf{P}^n_R$ by descending induction on $i$. Clearly the result
holds for $i \geq n + 1$ by (2). Suppose we know the result for
$i + 1$ and we want to show the result for $i$. (If $i = 0$, then
part (4) is vacuous.) Let $\mathcal{F}$ be a coherent sheaf on
$\mathbf{P}^n_R$. Choose a surjection as in (1) and denote
$\mathcal{G}$ the kernel so that we have a short exact sequence
$$
0 \to \mathcal{G} \to
\bigoplus\nolimits_{j = 1, \ldots, r}
\mathcal{O}_{\mathbf{P}^n_R}(d_j)
\to
\mathcal{F} \to 0
$$
By Lemma \ref{lemma-coherent-abelian-Noetherian}
we see that $\mathcal{G}$ is coherent. The long exact
cohomology sequence gives an exact sequence
$$
H^i(\mathbf{P}^n_R, \bigoplus\nolimits_{j = 1, \ldots, r}
\mathcal{O}_{\mathbf{P}^n_R}(d_j))
\to
H^i(\mathbf{P}^n_R, \mathcal{F})
\to
H^{i + 1}(\mathbf{P}^n_R, \mathcal{G}).
$$
By induction assumption the right $R$-module is finite and by
Lemma \ref{lemma-cohomology-projective-space-over-ring} the left
$R$-module is finite. Since $R$ is Noetherian it follows immediately
that $H^i(\mathbf{P}^n_R, \mathcal{F})$ is a finite $R$-module.
This proves the induction step for assertion (3).
Since $\mathcal{O}_{\mathbf{P}^n_R}(d)$ is invertible
we see that twisting on $\mathbf{P}^n_R$ is an exact functor (since
you get it by tensoring with an invertible sheaf, see
Constructions, Definition \ref{constructions-definition-twist}).
This means that for all $d \in \mathbf{Z}$ the sequence
$$
0 \to \mathcal{G}(d) \to
\bigoplus\nolimits_{j = 1, \ldots, r}
\mathcal{O}_{\mathbf{P}^n_R}(d_j + d)
\to
\mathcal{F}(d) \to 0
$$
is short exact. The resulting cohomology sequence is
$$
H^i(\mathbf{P}^n_R, \bigoplus\nolimits_{j = 1, \ldots, r}
\mathcal{O}_{\mathbf{P}^n_R}(d_j + d))
\to
H^i(\mathbf{P}^n_R, \mathcal{F}(d))
\to
H^{i + 1}(\mathbf{P}^n_R, \mathcal{G}(d)).
$$
By induction assumption we see the module on the right is zero
for $d \gg 0$ and by the computation in
Lemma \ref{lemma-cohomology-projective-space-over-ring}
the module on the left is zero as soon as $d \geq -\min\{d_j\}$
and $i \geq 1$. Hence the induction step for assertion (4).
This concludes the proof of (3) and (4).

\medskip\noindent
In order to prove (5) note that for all sufficiently large $d$
the map
$$
H^0(\mathbf{P}^n_R, \bigoplus\nolimits_{j = 1, \ldots, r}
\mathcal{O}_{\mathbf{P}^n_R}(d_j + d))
\to
H^0(\mathbf{P}^n_R, \mathcal{F}(d))
$$
is surjective by the vanishing of $H^1(\mathbf{P}^n_R, \mathcal{G}(d))$
we just proved. In other words, the module
$$
M_k
=
\bigoplus\nolimits_{d \geq k} H^0(\mathbf{P}^n_R, \mathcal{F}(d))
$$
is for $k$ large enough a quotient of the corresponding module
$$
N_k
=
\bigoplus\nolimits_{d \geq k} H^0(\mathbf{P}^n_R,
\bigoplus\nolimits_{j = 1, \ldots, r}
\mathcal{O}_{\mathbf{P}^n_R}(d_j + d)
)
$$
When $k$ is sufficiently small (e.g.\ $k < -d_j$ for all $j$) then
$$
N_k = \bigoplus\nolimits_{j = 1, \ldots, r}
R[T_0, \ldots, T_n](d_j)
$$
by our computations in Section \ref{section-cohomology-projective-space}.
In particular it is finitely generated.
Suppose $k \in \mathbf{Z}$ is arbitrary.
Choose $k_{-} \ll k \ll k_{+}$.
Consider the diagram
$$
\xymatrix{
N_{k_{-}} & N_{k_{+}} \ar[d] \ar[l] \\
M_k & M_{k_{+}} \ar[l]
}
$$
where the vertical arrow is the surjective map above and
the horizontal arrows are the obvious inclusion maps.
By what was said above we see that $N_{k_{-}}$ is a finitely
generated $R[T_0, \ldots, T_n]$-module. Hence $N_{k_{+}}$ is
a finitely generated $R[T_0, \ldots, T_n]$-module because it
is a submodule of a finitely generated module and the ring
$R[T_0, \ldots, T_n]$ is Noetherian. Since the vertical arrow
is surjective we conclude that $M_{k_{+}}$ is a finitely
generated $R[T_0, \ldots, T_n]$-module. The quotient
$M_k/M_{k_{+}}$ is finite as an $R$-module since it is a
finite direct sum of the finite $R$-modules
$H^0(\mathbf{P}^n_R, \mathcal{F}(d))$ for $k \leq d < k_{+}$.
Note that we use part (3) for $i = 0$ here. Hence
$M_k/M_{k_{+}}$ is a fortiori a finite $R[T_0, \ldots, T_n]$-module.
In other words, we have sandwiched $M_k$ between two finite
$R[T_0, \ldots, T_n]$-modules and we win.
\end{proof}